\documentclass[11pt]{article}
\usepackage{amsmath,amssymb,booktabs}
\usepackage[margin=1in]{geometry}

\title{Step 235: Cross-Track CRCFT Modes on Navier--Stokes Components}
\author{Six Birds Foundations III}
\date{2026-05-16}

\begin{document}
\maketitle

\section{Purpose}

This step tests whether the CRCFT modes classify the Navier--Stokes residual tree. The modes are:
\[
\text{TE},\qquad \text{CTMT},\qquad \text{BF}.
\]
In the NS track, CTMT adapts from matrix-element terminality to PDE-estimate terminality.

\section{NS residual structure}

The NS track is diagnostic-complete. The root target is the physical-native adequacy comparison
\[
\Xi(D^{\rm BG}\mid L^{\rm phys})\leq \Omega_{\rm BG}.
\]
The residual chain recorded in the NS map is
\[
\Xi(D^{\rm BG}\mid L^{\rm phys})
\leadsto
\Omega_{\rm src}
\leadsto
\Omega_{\rm gap}
\leadsto
\Omega_{\rm rem}
\leadsto
\Omega_{\rm bb}
\leadsto
\Omega_{\rm mm}
\leadsto
\Omega_{\rm amp}.
\]
The remaining external theorem is the sector-matched radius-window estimate
\[
\lambda_{J(t)}\rho(t)\geq \kappa\log \lambda_{J(t)}+h(t)
\]
on the fixed/exhaustive BG mismatch ledger, together with BKM/BG time integrability.

\section{Mode table}

\begin{center}
\begin{tabular}{p{0.32\linewidth}p{0.16\linewidth}p{0.42\linewidth}}
\toprule
Component & Mode & Reason\\
\midrule
EXT1 radius-window product theorem & TE/CTMT & Target-strength final theorem; proof is terminal PDE estimate.\\
EXT2 BKM/BG time gate & TE & Continuation time class is target-strength.\\
EXT3 sector match & BF & Candidate sparse/Gevrey sector does not transport unless it matches the BG sector.\\
EXT4 tail/lift budget & CTMT & Quantitative tail and tangent-lift estimates are terminal.\\
\(\Omega_{\rm src}\) & BF/CTMT & Raw nonlinear shadows need lawful source-admission decomposition.\\
\(\Omega_{\rm gap},\Omega_{\rm rem}\) & CTMT & Localization and remainder budgets are terminal PDE components.\\
\(\Omega_{\rm bb}\) & BF/CTMT & Packing-only transport fails unless amplitude threshold is paid.\\
\(\Omega_{\rm mm}\) & BF & It is the sector-mismatch bridge failure.\\
\(\Omega_{\rm amp}\) & TE/CTMT & Final sharp obstruction; closure requires EXT1.\\
\bottomrule
\end{tabular}
\end{center}

\section{Literature audit}

Beale--Kato--Majda gives the vorticity continuation criterion. Caffarelli--Kohn--Nirenberg gives partial regularity. Constantin--Fefferman supplies geometric depletion criteria. Koch--Tataru gives critical-space well-posedness. Tao's averaged Navier--Stokes blowup and Buckmaster--Vicol weak-solution nonuniqueness provide guardrails. Bradshaw--Grujic and Grujic--Xu provide frequency/sparseness route templates. None supplies the exact EXT1 theorem on the fixed BG mismatch ledger.

\section{Verdict}

\[
\boxed{\texttt{V\_ns\_CRCFT\_modes\_verified}}
\]

CRCFT is now verified on four track instances: RH, BSD, Hodge, and NS. No Navier--Stokes proof is claimed.

\end{document}
